% ============================================================================
% Chapter 8: Windows Desktop Environment
% ============================================================================

\chapter{Windows Desktop Environment}
\label{ch:desktop}

\begin{objectives}
  \item Identify the main elements of the Windows desktop
  \item Explain the purpose of the Taskbar, Start Menu, and System Tray
  \item Use desktop icons and shortcuts effectively
  \item Navigate the Windows desktop with confidence
\end{objectives}

\bytetip[fun]{The desktop is your computer's ``home screen'' --- like the home screen on a phone, but with WAY more power. Let's explore every corner of it!}

% ---------------------------------------------------------------
\section{The Windows Desktop}
\label{sec:windows-desktop}
\index{desktop}\index{Windows!desktop}

When Windows finishes loading, you see the \hl{desktop}. This is your main workspace --- everything starts here.

% --- Figure 8.1: Annotated Desktop ---
\begin{figure}[H]
\centering
\begin{tikzpicture}[scale=0.75]
  % Screen outline
  \draw[NavyBlue, line width=2pt, rounded corners=4pt] (0,0) rectangle (16,10);
  
  % Desktop background
  \fill[NavyBlue!5] (0.15,0.8) rectangle (15.85,9.85);
  
  % Taskbar
  \fill[NavyBlue!80] (0.15,0) rectangle (15.85,0.8);
  
  % Start button
  \fill[ModuleBlue] (0.15,0) rectangle (1.0,0.8);
  \node[text=white, font=\scriptsize\bfseries] at (0.575,0.4) {\faWindows};
  
  % Pinned taskbar icons
  \foreach \x/\icon in {1.5/\faFolder, 2.2/\faChrome, 2.9/\faFileWord, 3.6/\faStore}{
    \node[text=white!80, font=\scriptsize] at (\x, 0.4) {\icon};
  }
  
  % System tray (right side)
  \node[text=white!70, font=\tiny\sffamily] at (13.5,0.4) {\faWifi\quad \faVolumeUp\quad \faBatteryFull};
  \node[text=white!90, font=\tiny\sffamily] at (15.0,0.4) {12:30 PM};
  
  % Desktop icons
  \foreach \y/\icon/\lbl in {8.5/\faTrash/Recycle Bin, 7.0/\faChrome/Chrome, 5.5/\faFolder/Documents}{
    \node[text=NavyBlue, font=\large] at (1.2, \y) {\icon};
    \node[font=\tiny\sffamily, text=NavyBlue] at (1.2, \y-0.5) {\lbl};
  }
  
  % Open window (File Explorer)
  \fill[white] (4,2.5) rectangle (13,8.5);
  \fill[ModuleBlue!80] (4,7.8) rectangle (13,8.5);
  \node[text=white, font=\tiny\sffamily\bfseries, anchor=west] at (4.3,8.15) {File Explorer};
  \node[text=white, font=\tiny\bfseries, anchor=east] at (12.7,8.15) {$\times$};
  
  % Labels with arrows (outside screen)
  % Desktop
  \draw[FunPurple, line width=1.2pt, ->] (-2.5, 6) -- (2, 6);
  \node[left, font=\scriptsize\sffamily\bfseries, text=FunPurple, align=right] at (-2.5,6) {Desktop\\Icons};
  
  % Taskbar
  \draw[FunCoral, line width=1.2pt, ->] (-2.5, 0.4) -- (0.1, 0.4);
  \node[left, font=\scriptsize\sffamily\bfseries, text=FunCoral, align=right] at (-2.5,0.4) {Taskbar};
  
  % Start
  \draw[LabGreen, line width=1.2pt, ->] (0.575, -1.2) -- (0.575, -0.1);
  \node[below, font=\scriptsize\sffamily\bfseries, text=LabGreen] at (0.575,-1.2) {Start Button};
  
  % System tray
  \draw[NoteOrange, line width=1.2pt, ->] (17.5, 0.4) -- (15.9, 0.4);
  \node[right, font=\scriptsize\sffamily\bfseries, text=NoteOrange, align=left] at (17.5,0.4) {System Tray\\(Notifications)};
  
  % Window
  \draw[ModuleBlue, line width=1.2pt, ->] (17, 7) -- (13.1, 7);
  \node[right, font=\scriptsize\sffamily\bfseries, text=ModuleBlue, align=left] at (17,7) {Open\\Window};
\end{tikzpicture}
\caption{The Windows desktop with all main elements labelled}
\label{fig:desktop-annotated}
\end{figure}

% --- Table 8.1: Desktop Elements Quick Reference ---
\begin{table}[H]
\centering
\caption{Desktop elements quick reference}
\label{tab:desktop-elements}
\renewcommand{\arraystretch}{1.3}
\begin{tabularx}{\textwidth}{c l X}
\toprule
\rowcolor{HeaderRow}
\textcolor{white}{\textbf{Icon}} & \textcolor{white}{\textbf{Element}} & \textcolor{white}{\textbf{What It Does}} \\
\midrule
\rowcolor{RowLight}
\faDesktop & Desktop & Main working area where icons and windows appear \\
\rowcolor{RowDark}
\faMinus & Taskbar & Shows open programs and pinned apps at the bottom \\
\rowcolor{RowLight}
\faWindows & Start Menu & Opens apps, settings, power options, and search \\
\rowcolor{RowDark}
\faFolder & Icons & Shortcuts to files, folders, and programs \\
\rowcolor{RowLight}
\faBell & Notification Area & Alerts, date/time, Wi-Fi, volume, and battery \\
\rowcolor{RowDark}
\faTrash & Recycle Bin & Deleted files go here (can be recovered!) \\
\bottomrule
\end{tabularx}
\end{table}

\begin{notebox}
The \textbf{Recycle Bin} is like a rubbish bin for your computer. When you delete a file, it goes to the Recycle Bin first. You can restore it if you change your mind! Files are only permanently deleted when you \textbf{empty} the Recycle Bin.
\end{notebox}

\bytetip[tip]{You can customise your desktop! Right-click on an empty area of the desktop and try ``Personalise'' to change the wallpaper, colours, and theme. Make your computer feel like \emph{yours}!}

\begin{funfact}
The Recycle Bin icon on your desktop actually changes its appearance! When it's empty, it shows an empty bin. When it contains deleted files, it shows a bin with crumpled paper inside. Nice detail, right?
\end{funfact}

\begin{reviewqs}
  \item Name five elements of the Windows desktop.
  \item What is the purpose of the Taskbar?
  \item How do you access the Start Menu?
  \item What is the Recycle Bin, and where do deleted files go?
  \item What information can you find in the System Tray (notification area)?
\end{reviewqs}

\begin{challenge}
Customise your desktop: change the wallpaper, rearrange the icons, and pin your favourite app to the taskbar. Then take a screenshot (\key{PrtSc} key) and paste it into MS Paint. Save it as ``MyDesktop.png.''\end{challenge}
